% Task B 报告骨架 —— 最终与 Part A 合并成一份 PDF（任务书 §4.1、§4.4）。
% 报告用英文写；结构照 §4.1 的 B.1–B.6。
% 文件名保持 main.tex；正文你自己写，这里只留提纲。

% B.1 Overview & Constraints
%     约束清单（禁用框架）、自己立的纯度守卫测试、环境与版本留痕
%     （sys.executable / numpy / torch 版本 + 说明实验跑在哪个解释器上）

% B.2 Implementation Discipline          ← 加分项：你怎么保证不越界
%     目录 / 环境隔离方案（task_b 子项目、npnn 包名、数据只读共享、不装包）
%     + 纯度守卫测试的机制与**输出**（证明而不是声称）

% B.3 Module Implementations (Scoring Item 1)              ★ 40 分
%     B.3.1 Linear
%     B.3.2 GELU
%     B.3.3 CrossEntropyLoss
%     B.3.4 Adam
%     B.3.5 Conv2d（im2col 说明、朴素 oracle、重叠累加的处理、耗时对比）
%     B.3.6 SwitchNorm2d（三种退化测试、λ 梯度、权重演化曲线）
%     B.3.7 Verification summary table（每模块 rel_err 汇总成一张表）
%     每个模块按同一模板：公式 → 推导要点 → 对拍结果表

% B.4 End-to-End Training & Comparison with PyTorch (Scoring Item 2)   ★ 10 分
%     B.4.1 网络设计与预算取舍（引用 §0.5 的计时外推，解释砍了什么、为什么）
%     B.4.2 整网对拍协议与结果（L2 核心证据：loss 差值曲线，对数轴）
%     B.4.3 端到端精度与速度对比（精度是否相近 / 轨迹是否一致 / 慢多少倍）
%     B.4.4 性能工程（im2col、线程、dtype、连续性）

% B.5 What I Learned        ← 诚实 + 具体 + 带数字，3–5 个点，失败与返工是加分项
% B.6 References

% 图表清单见 §4.2（9 项，照着做不会漏）。
% ⚠️ 所有误差都给量级（max|Δ| = 3.2e-11 这种），不要写「基本一致」。
% ⚠️ 图全部从 task_b/runs/ 的落盘数据画（notebook 13），不要现训现画。
